% ---------------------------------------------------------------------------
% Ch. 8 -- Selection of the DER Q(V) dead-zone half-width.
%
% Supersedes docs/ch8_deadband_window_selection.tex, whose results were produced
% with the rev-1 sensitivity reduction (see Sec. 3).
%
% To \input into the thesis, delete from \documentclass down to
% \tableofcontents, and the final \end{document}.
%
% Figures: results/deadband_selection/figures/fig_*.pdf
% Data:    results/deadband_selection/deadband_{metrics,activity,threshold,
%          disturbance}.csv
% ---------------------------------------------------------------------------
\documentclass[11pt,a4paper]{article}

\usepackage[utf8]{inputenc}
\usepackage[T1]{fontenc}
\usepackage{amsmath,amssymb}
\usepackage{booktabs}
\usepackage{siunitx}
\usepackage{graphicx}
\usepackage[margin=2.5cm]{geometry}
\usepackage[hidelinks]{hyperref}

\sisetup{detect-weight=true,detect-family=true}
\graphicspath{{../results/deadband_selection/figures/}%
              {../results/deadband_n1/figures/}}

\newcommand{\pu}{\ensuremath{\mathrm{p.u.}}}
\newcommand{\Mvar}{\ensuremath{\mathrm{Mvar}}}
\newcommand{\qss}{\textsc{qss}}
\newcommand{\rms}{\textsc{rms}}
\newcommand{\dz}{\ensuremath{\delta}}

\title{Selection of the DER $Q(V)$ dead-zone half-width}
\author{}
\date{2026--08--02}

\begin{document}
\maketitle
\tableofcontents

%======================================================================
\section{The parameter and the question}

The distribution-connected DER follow a local reactive-power droop $Q(V)$ with a
dead zone of half-width $\dz$ about a re-anchored voltage reference: the droop
is inactive while $|V-V_\mathrm{anchor}| < \dz$ and engages beyond it. $\dz$ is
therefore an \emph{amplitude threshold on the controller's input}, and it is a
property of the characteristic rather than of any operating point.

Two effects are asserted to bound it. Too narrow, and the droop answers every
small variation: the DER move continuously and the outer online feedback
optimisation (OFO) repeatedly re-anchors against them. Too wide, and the droop
is inactive between dispatches, so no local support acts and the interface
reactive power drifts until the next OFO step.

This chapter reports three experiments. The first --- sweeping $\dz$ against
annual profile windows --- is the obvious design and \textbf{it fails}; why it
fails is itself a result (\S\ref{sec:sweep}). The second measures the mechanism
the dead zone governs rather than its consequence (\S\ref{sec:activity}). The
third measures the threshold directly by controlled excitation
(\S\ref{sec:threshold}) and yields the design rule of \S\ref{sec:rule}.

\subsection{Plant, actuators, controlled outputs}

\begin{itemize}
  \item \textbf{Plant.} Modified IEEE~39-bus transmission system, scenario
    \texttt{rural\_700}, with four \SI{110}{\kilo\volt} distribution underlays.
    $\mathrm{DSO}_3$ carries $2\times$ installed DER and $2\times$ active-load
    base; its reactive load is not scaled. Every DSO has zero constant reactive
    load and follows $Q_\mathrm{load}(t)=\SI{500}{\Mvar}\cdot m_\mathrm{qload}(t)$.
  \item \textbf{Actuators.} DER reactive power (continuous, $Q(V)$ droop with
    dead zone), coupler three-winding OLTCs, machine two-winding OLTCs, MSC/MSR
    switched shunts.
  \item \textbf{Controlled outputs.} Reactive power at the EHV--HV interfaces,
    and nodal voltages within their band.
  \item \textbf{DER capability.} The physical VDE-AR-N-4120 diagrams throughout.
    Reactive capability is contingent on active power, which matters for window
    selection: a park at $P=0$ has none (\S\ref{sec:degenerate}).
  \item \textbf{Controllers see only cached sensitivities} and their own
    measurements. No controller is given the plant equations.
\end{itemize}

\subsection{Excitation}
\label{sec:excitation}

Two disturbance classes are used, and \textbf{no contingency is applied
anywhere}: every run records \texttt{contingencies = []},
\texttt{measurement\_noise.enabled = False} and
\texttt{enable\_reachability\_guard = False}.

\begin{description}
  \item[Profiled operation.] The annual load/DER time series. This is the
    small-signal regime the dead zone is meant to filter.
  \item[Exogenous load step.] A step multiplying every active-load profile from
    a fixed instant. Applied to the \emph{interpolated} profile frame, so it is
    a true step at the \SI{20}{\second} dispatch resolution --- perturbing the
    \SI{15}{\minute} source would have been smeared into a ramp by
    interpolation. It reaches both plants through paths they already support
    and is \emph{not} a contingency: no element is switched.
\end{description}

The absence of measurement noise bounds every conclusion below. A dead zone's
textbook justification is rejection of noise on the measured voltage; that
mechanism is absent by construction, so the narrow-side penalty reported here is
carried by genuine profile variation alone. Adding noise should penalise narrow
dead bands further, making every $\dz^\star$ reported here a \emph{lower} bound
on the value appropriate to a noisy plant. The direction is argued, not measured.

%======================================================================
\section{Provenance: the sensitivity reduction was corrected}
\label{sec:rev2}

Each controller acts on cached sensitivities derived from a Ward-reduced
network. On 2026-08-01 that reduction was found to be wrong in three
independent ways, and every result in this chapter postdates the correction.

Verified with \texttt{tools/check\_reduction\_fidelity.py}, which asks whether
each reduced network, solved alone, reproduces the combined solution at the
buses it retains --- the property the reduction claims. Convergence is
\emph{not} that test: a reduced net with stranded boundary flows converges
happily to a different operating point.

\begin{table}[htbp]
\centering
\begin{tabular}{lll}
\toprule
Defect & Effect & Fix \\
\midrule
DSO reductions solved on the wrong power-flow branch
  & \SIrange{0.10}{0.36}{\pu} from the combined solution; tertiary buses
    collapsed to \SI{0}{\pu}
  & DC start before flat \\
Sub-network's own auxiliary buses dropped
  & \SI{132}{\mega\watt} of the DSO's own injection missing
  & keep \texttt{internal\_aux} buses \\
Non-slack boundary couplers pinned at $P=0$
  & multi-coupler DSOs forced their whole real exchange through one transformer
  & carry cached $p_\mathrm{hv}$ \\
TSO zones used generator \emph{setpoints}
  & \SI{180}{\mega\watt} phantom injection under distributed slack
  & use cached \texttt{res\_gen} \\
\bottomrule
\end{tabular}
\caption{Defects in the reduced-network sensitivities, corrected 2026-08-01.
All four inputs are locally measurable: PCC flows, own topology, own machine
telemetry.}
\end{table}

After correction all sixteen DSO reductions (4 DSOs $\times$ 4 windows)
reproduce the combined solution to \textbf{0.000000\,\pu}; the TSO zones improved
up to $5\times$ and retain a residual of \SIrange{0.012}{0.018}{\pu}, which is a
linearisation-point offset in a smooth regime and is accepted.

The correction is not cosmetic: at 2016-01-05 08:00, $\dz=0.005$, the
interface-$Q$ metric moved \SI{0.4410}{\Mvar} $\rightarrow$ \SI{0.4929}{\Mvar}
($\approx\!12\,\%$). Runs record \texttt{sensitivity\_reduction\_rev}, and the
admission filter requires rev~2, so pre-correction runs cannot enter these
curves.

%======================================================================
\section{Windows}
\label{sec:windows}

Five live windows spanning net infeed $-117$ to $+\SI{2200}{\mega\watt}$, plus
one retained as a null result.

\begin{table}[htbp]
\centering
\begin{tabular}{lS[table-format=4.1]S[table-format=4.1]S[table-format=+5.1]c}
\toprule
Window & {DER $P$} & {Load $P$} & {Net infeed} & Parks with \\
 & {[\si{\mega\watt}]} & {[\si{\mega\watt}]} & {[\si{\mega\watt}]}
 & zero $Q$ cap. \\
\midrule
2016-02-22 13:00 & 2262.4 & 2379.9 & -117.4 & 16/44 \\
2016-01-05 08:00 & 2324.5 & 1916.0 &  408.6 & 16/44 \\
2016-01-15 03:00 & 2605.7 & 1800.7 &  805.0 & 16/44 \\
2016-12-18 14:00 & 2741.8 & 1374.7 & 1367.1 & 16/44 \\
2016-05-01 16:00 & 3458.6 & 1258.6 & 2200.0 & {--} \\
\addlinespace
2016-07-15 03:00 &   29.2 & 1055.3 & -1026.1 & \textbf{44/44} \\
\bottomrule
\end{tabular}
\caption{Windows characterised by current-topology quantities. The last is
degenerate (\S\ref{sec:degenerate}).}
\label{tab:windows}
\end{table}

Three points of method:

\paragraph{The screening figures that selected the windows are void.} Three
windows were drawn from a Tier-1 seasonal screening measured on an
\emph{earlier topology}. That ranking does not transfer: the window it called
the annual maximum turned out to be the least demanding of the set. The figures
are recorded as provenance only and are never quoted as properties of these
runs.

\paragraph{A deep-import window with capable DER does not exist in this data.}
DER output and load are positively correlated, so the hours with reactive
capability are the high-generation hours. Screening all $35\,136$ quarter-hours,
\emph{no} hour combines $P_\mathrm{DER}\geq\SI{2200}{\mega\watt}$ with net
infeed below $-\SI{300}{\mega\watt}$. The import side of the axis therefore
stops at $-\SI{117}{\mega\watt}$, and any statement about the voltage-drop
regime rests on that single window.

\paragraph{Feasibility is empirical.} 2016-05-01 16:00 aborted with a
non-convergent power flow under rev~1 and converges under rev~2, because
removing the phantom generation of \S\ref{sec:rev2} gave its reduced net a
solution. Offline screening did not predict this in either direction.

\subsection{A degenerate window}
\label{sec:degenerate}

At 2016-07-15 03:00 --- July, 03:00, no photovoltaic infeed and a multi-day wind
lull --- aggregate DER infeed is \SI{29.2}{\mega\watt}. Under the VDE diagram
\emph{all} 44 parks report zero reactive capability and both zone PCC capability
intervals collapse to $[0,0]\,\si{\Mvar}$. With no DER able to move, the $Q(V)$
characteristic never binds: two dead bands an order of magnitude apart return
identical results to six digits.

This is a property of the operating point, not a defect, and it bounds the
design rule: \textbf{dead-band selection presupposes DER reactive headroom.}
The window is excluded from every optimum and retained as a null result. It also
serves as a determinism control --- because $\dz$ provably cannot act there,
its identical outputs prove the pipeline is deterministic, which matters for
\S\ref{sec:sweep}.

%======================================================================
\section{Experiment 1: the profiled sweep, and why it fails}
\label{sec:sweep}

Full factorial, $\dz \in \{0, 0.001, 0.0025, 0.005, 0.0075, 0.01, 0.015, 0.02,
0.03\}\,\pu$ crossed with the five live windows: 45 closed-loop
co-simulations against the PowerFactory \rms{} plant, each with a quasi-static
reference leg.

\begin{figure}[htbp]
\centering
\includegraphics[width=0.78\textwidth]{fig_tracking.pdf}
\caption{Interface-$Q$ tracking error against dead-zone half-width, one curve
per window. Only the mildest window ($+\SI{409}{\mega\watt}$) gives a clean
interior minimum; the stressed windows are non-monotone.}
\label{fig:tracking}
\end{figure}

\begin{figure}[htbp]
\centering
\includegraphics[width=0.78\textwidth]{fig_dsvoltage.pdf}
\caption{DS group voltage deviation against $\dz$. Every exporting window is
best at $\dz=0$; the importing window ($-\SI{117}{\mega\watt}$) has an interior
optimum, i.e.\ the trade-off reverses sign with the direction of interface flow.}
\label{fig:dsv}
\end{figure}

\subsection{Result: no optimum can be defined}

\begin{itemize}
  \item \textbf{No $\dz$ is Pareto-optimal in every window} over
    (interface $Q$, DS voltage) --- the intersection is empty
    (Fig.~\ref{fig:tradeoff}).
  \item The interface-$Q$ argmin scatters over $0.0025 \dots 0.03$ with no
    relation to loading: $\mathrm{CV}=0.715$ across the five windows.
  \item The stressed windows' curves are non-monotone with \SIrange{20}{25}
    {\percent} scatter.
\end{itemize}

\begin{figure}[htbp]
\centering
\includegraphics[width=0.78\textwidth]{fig_tradeoff.pdf}
\caption{The two controlled quantities against each other, with $\dz$ as the
path parameter. A monotone trade-off would trace a simple curve; the stressed
windows fold back on themselves.}
\label{fig:tradeoff}
\end{figure}

\subsection{Why: the metric is not a function of $\dz$}

The scatter is \emph{structural}, not statistical. The degenerate window
(\S\ref{sec:degenerate}) reproduces to six digits, so the pipeline is
deterministic; different $\dz$ therefore reach genuinely different equilibria.
This is the dead-zone multi-equilibrium behaviour --- the same non-uniqueness
identified between the quasi-static and \rms{} representations --- here
dominating the stressed operating points rather than appearing as isolated
anomalies.

Four design flaws compound it: one realisation per cell, so a single trajectory
samples one basin arbitrarily; a $\dz$-dependent initial condition, because the
$Q(V)$ seeding uses the dead band; a metric mixing transient with steady state;
and no measurement of the mechanism the dead band actually governs.

\textbf{Taking an argmin over a bistable metric is not a valid selection
procedure}, and the remainder of this chapter does not attempt it.

Two findings from this experiment do survive, because they do not depend on
locating a minimum:

\begin{enumerate}
  \item \textbf{$\dz$ must be non-zero.} At $\dz=0$ the interface-$Q$ error is
    roughly $3\times$ its best value in every window. The collapse is abrupt:
    $\dz=\SI{0.001}{\pu}$ already recovers about \SI{90}{\percent} of it.
  \item \textbf{The voltage trade-off reverses with the direction of interface
    flow} (Fig.~\ref{fig:dsv}). Exporting windows want $\dz=0$ for DS voltage;
    the importing window has an interior optimum at \SI{0.0075}{\pu}.
\end{enumerate}

%======================================================================
\section{Experiment 2: measuring the mechanism}
\label{sec:activity}

The dead zone governs how much the local droop \emph{moves}. Counting actuator
motion rather than reading a possibly-bistable end state gives a metric that
cannot be confounded by which equilibrium a run settles into.

\begin{figure}[htbp]
\centering
\includegraphics[width=0.78\textwidth]{fig_chatter.pdf}
\caption{Direction reversals of the DER reactive power, per park per dispatch
interval, against $\dz$. Monotone with a clear knee in every window --- compare
the same runs in Fig.~\ref{fig:tracking}.}
\label{fig:chatter}
\end{figure}

Defining the knee as the smallest $\dz$ within \SI{10}{\percent} of that
window's floor:

\begin{center}
\begin{tabular}{lcc}
\toprule
 & mean & CV \\
\midrule
chatter knee & \SI{0.0100}{\pu} & \textbf{0.418} \\
interface-$Q$ argmin & \SI{0.016}{\pu} & 0.715 \\
\bottomrule
\end{tabular}
\end{center}

The chatter knee is about twice as well determined, and it measures the quantity
$\dz$ directly controls. It is still a $3\times$ range rather than a number,
which motivates \S\ref{sec:threshold}.

\paragraph{A normalisation hypothesis, tested and refuted.} If $\dz$ were simply
a threshold relative to local voltage variability, normalising by the measured
$\sigma_V$ (which ranges \SIrange{0.00128}{0.00481}{\pu} across the windows)
should collapse the curves. It does not: the knee CV goes $0.418 \rightarrow
0.427$ and the argmin CV $0.715 \rightarrow 0.747$. Both marginally
\emph{worse}. $\dz$ is not a constant multiple of $\sigma_V$.

%======================================================================
\section{Experiment 3: the threshold, measured directly}
\label{sec:threshold}

Rather than infer the threshold from \SI{300}{\second} of uncontrolled
multi-frequency excitation, apply one disturbance of known amplitude and find
the $\dz$ at which the droop stops responding. One operating point
(2016-01-05 08:00 --- the only window whose response is unambiguous, and the
lowest $\sigma_V$, so the step dominates the background).

\begin{figure}[htbp]
\centering
\includegraphics[width=0.78\textwidth]{fig_threshold.pdf}
\caption{Post-step DER traverse against $\dz$ at two excitation amplitudes.
Dotted verticals mark each amplitude's open-loop $|\Delta V|$. The response
collapses to the quiet floor as $\dz$ crosses the excitation.}
\label{fig:threshold}
\end{figure}

\begin{table}[htbp]
\centering
\begin{tabular}{lcccccc}
\toprule
$\dz$ [\pu] & 0.0025 & 0.005 & 0.0075 & 0.01 & 0.015 & 0.02 \\
\midrule
$\times1.10$, $|\Delta V| = \SI{0.00255}{\pu}$
  & 0.113 & \textbf{0.049} & 0.051 & 0.054 & 0.055 & 0.055 \\
$\times1.25$, $|\Delta V| = \SI{0.00690}{\pu}$
  & 0.752 & 0.476 & 0.161 & \textbf{0.081} & 0.058 & 0.058 \\
\bottomrule
\end{tabular}
\caption{Post-step DER traverse [\si{\Mvar} per park per interval]. Bold marks
the first $\dz$ at which the response has essentially reached the quiet floor.
In both rows that happens as $\dz$ passes the excitation amplitude.}
\label{tab:threshold}
\end{table}

At $|\Delta V| = \SI{0.00255}{\pu}$ the response falls to the floor between
$\dz = 0.0025$ and $0.005$; at $|\Delta V| = \SI{0.00690}{\pu}$ between
$\dz=0.005$ and $0.0075$. \textbf{In both cases the transition brackets the
excitation amplitude}, so the engagement threshold satisfies
\begin{equation}
  \dz_\mathrm{engage} \;\approx\; |\Delta V|,
  \qquad k = \dz_\mathrm{engage}/|\Delta V| \approx 1 ,
  \label{eq:threshold}
\end{equation}
at two independent amplitudes differing by a factor of~$2.7$. The dead zone
behaves as specified: the droop acts when the deviation exceeds $\dz$ and is
silent when it does not.

\paragraph{The excitation is not independent of $\dz$.} Measured $|\Delta V|$ at
$\times1.25$ runs from \SI{0.00153}{\pu} at $\dz=0$ to \SI{0.00690}{\pu} at
$\dz\geq0.015$: a tight dead band makes the droop act, and acting suppresses the
very deviation it responds to. The open-loop excitation must therefore be read
at wide $\dz$, where the droop is inactive. \emph{The dead band partly determines
its own input}, which is a closed-loop property worth stating explicitly.

\paragraph{The threshold is measured on the TS parks.} Equation
\eqref{eq:threshold} was obtained from the aggregate traverse, which the
TS-connected parks dominate by a factor of \numrange{10}{40}; it is therefore in
effect a TS result, and it does not by itself license one $\dz$ for both
voltage levels. Resolving the excitation per park at wide $\dz$ (open loop,
$\times1.25$) gives

\begin{center}
\begin{tabular}{lcccc}
\toprule
 & $n$ & median $|\Delta V|$ [\pu] & max [\pu] & max/median \\
\midrule
TS parks & 4  & 0.00645 & 0.00813 & 1.26 \\
DS parks & 40 & 0.00673 & 0.00907 & 1.35 \\
\bottomrule
\end{tabular}
\end{center}

\noindent The two populations differ in median by under \SI{4}{\percent}, and
neither spreads by more than $1.35\times$ about its own median. Referenced to
the \emph{maximum} park deviation rather than the mean --- the appropriate
reference, since the dead band must be chosen against the worst park it has to
keep silent --- $k \approx 1$ holds for both levels. A single uniform $\dz$ is
therefore justified \emph{for this disturbance class}; whether it survives a
disturbance that loads one level preferentially is a separate question, taken up
in \S\ref{sec:rule}.

\paragraph{Large dead bands are unreachable by \emph{uniform} disturbances.} A
static scan gives the load step needed to induce a given deviation: under a
uniform increase at all profiled buses, exceeding \SI{0.02}{\pu} requires roughly
\emph{tripling} system load ($4938 \rightarrow \SI{8770}{\mega\watt}$). That is a
stress test, not an event.

This conclusion is conditional on the disturbance being uniform, and that
qualification matters: a load step \emph{concentrated} at a single weak bus is
far more efficient. At the most severe distribution bus, \SI{43}{\mega\watt}
suffices for \SI{0.02}{\pu} --- some \num{90} times less energy than the uniform
step, with a spatial spread of \SI{0.046}{\pu} across the parks. A localised
event of entirely ordinary size therefore reaches deviations that the uniform
scan places beyond reach. The claim that should be carried forward is the
narrower one: \emph{under uniform loading} at this operating point,
$\dz \gtrsim 0.01$ leaves the local droop effectively disabled. Amplitudes at and beyond $\times1.5$ were attempted and
proved unusable for a different reason: they destabilise the cascade
(traverse \num{80.7} against a floor of \num{0.057}, and a settling-gate
failure), so they measure instability rather than threshold.

%======================================================================
\section{The design rule}
\label{sec:rule}

Equation~\eqref{eq:threshold} converts $\dz$ from a tuning parameter into a
stated design choice:

\begin{quote}
\textbf{Choose $\dz$ equal to the voltage deviation the local layer should
ignore. The droop then engages for anything larger and is silent for anything
smaller.}
\end{quote}

The performance results then bound the admissible range rather than selecting
within it:

\begin{center}
\begin{tabular}{lll}
\toprule
Bound & From & Value \\
\midrule
lower & $\dz=0$ costs $3\times$ in interface-$Q$ tracking & $\dz > 0$ \\
lower & chatter knee (\S\ref{sec:activity}) & $\dz \gtrsim 0.005$ \\
upper & DS voltage degrades monotonically (exporting) & as small as tolerable \\
upper & droop never engages on realistic events & $\dz \lesssim 0.01$ \\
\bottomrule
\end{tabular}
\end{center}

\subsection{Worked example: a \SI{0.01}{\pu} limit}

Suppose the requirement is that the local layer must act before the deviation
reaches \SI{0.01}{\pu}. By \eqref{eq:threshold}, engagement occurs when the
deviation exceeds $\dz$, so the requirement is
\[
  \dz \;\leq\; \SI{0.01}{\pu}.
\]
$\dz=\SI{0.01}{\pu}$ is the largest admissible value and satisfies the
requirement exactly at the limit. It also sits at the chatter knee, so it
captures essentially all of the achievable chatter reduction.

It is, however, the \emph{worst} admissible choice on the other two measures. At
the reference window, relative to their best values, $\dz=\SI{0.01}{\pu}$ costs
\SI{18}{\percent} in interface-$Q$ tracking (\num{0.580} against \num{0.493}
\si{\Mvar}) and \SI{73}{\percent} in step rejection (peak excess \num{0.335}
against \num{0.194} \si{\Mvar}), and it gives up a factor of $3.6$ in DS voltage
against $\dz=0$.

\textbf{Recommendation: $\dz = \SI{0.005}{\pu}$.} It satisfies the
\SI{0.01}{\pu} requirement with a factor-of-two margin, and it is where three
independent criteria at the reference window already coincide:

\begin{center}
\begin{tabular}{ll}
\toprule
Criterion & $\dz$ \\
\midrule
profiled interface-$Q$ optimum & 0.005 \\
step-rejection optimum & 0.005 \\
chatter knee & 0.0075 \\
\bottomrule
\end{tabular}
\end{center}

%======================================================================
\section{The dead band as a detector threshold: N-1 evidence}
\label{sec:n1}

\begin{quote}\small
\textbf{Provisional.} Measured on THREE operating windows spanning net infeed
\SIrange{-117}{+1367}{\mega\watt} (2016-02-22 13:00, 2016-01-05 08:00,
2016-12-18 14:00), 57 RMS co-simulations run 2026-08-02/04. One cell
($\dz=0.1$, gen~1, \SI{+409}{\mega\watt}) lost its traces to a failed result
export and is excluded.

\textbf{Supersedes the single-window draft.} That draft concluded that
$\dz=\SI{0.005}{\pu}$ was confirmed. Replication \emph{refutes} it: see
\S\ref{sec:n1choice}.
\end{quote}

\subsection{Why the earlier sections could not find an optimum}

The Q(V) layer exists to reject \emph{events}. Profile-driven voltage shifts are
the OFO's business: it re-optimises every \SI{180}{\second} (TSO) or
\SI{20}{\second} (DSO) and re-anchors each park's $V_\mathrm{ref}$ when it writes
a setpoint. The droop compares $|V-V_\mathrm{anchor}|$ against $\dz$, so $\dz$
does not discriminate voltage \emph{levels} --- it discriminates \emph{drift
since that park's last dispatch}.

That makes $\dz$ a \textbf{detector threshold}, and a detector has no optimum. It
has an admissible interval, bounded below by the disturbance it must ignore and
above by the one it must catch:

\begin{equation}
  \underbrace{\max_{\text{profile}} |\Delta V|}_{\text{OFO's business}}
  \;\lesssim\; \dz \;\lesssim\;
  \underbrace{\min_{\text{event}} |\Delta V|}_{\text{droop's business}} .
  \label{eq:detector}
\end{equation}

This is why $\dz^\star$ scattered with $\mathrm{CV}=0.715$ across windows in
\S\ref{sec:sweep}: the argmin was being fitted to a curve whose flat region
\emph{is} the answer.

\subsection{Method}

The disturbance is a synchronous-machine trip --- a credible N-1 --- delivered to
the RMS plant as an \texttt{EvtOutage} on the machine and its unit transformer.
Horizon \SI{600}{\second}; the trip fires at \SI{200}{\second}, i.e.
\SI{20}{\second} \emph{after} a TSO dispatch, which maximises the droop-only
exposure to \SI{160}{\second}. Each of the eight dead bands is run three times:
undisturbed twin, mild trip, severe trip. Every deviation is referenced to the
same-$\dz$ twin; the pre-trip deviation between a run and its twin is
\textbf{identically zero}, so the referencing contributes no artefact.

\paragraph{The disturbance is a transient, and that is the point.} The excursion
is not a sustained offset. At the worst park, with the droop disabled, a gen-7
trip gives

\begin{center}
\begin{tabular}{lccccccc}
\toprule
$t$ [\si{\second}] & 179 & 190 & 195 & 220 & 260 & 300 & 600 \\
\midrule
$V$ [\pu]          & 1.0441 & 0.9132 & 0.8949 & 0.9024 & 1.0052 & 1.0389 & 1.0559 \\
$|\Delta V|$ [\pu] & 0.0000 & 0.1315 & 0.1497 & 0.1424 & 0.0397 & 0.0060 & 0.0093 \\
\bottomrule
\end{tabular}
\end{center}

\noindent (the twin holds \SIrange{1.0441}{1.0466}{\pu} throughout, so the
deviation is attributable to the trip alone; pre-trip $|\Delta V|$ is
identically zero). The peak of \SI{0.2240}{\pu} occurs \SI{10.5}{\second} after
the trip; $|\Delta V|$ exceeds \SI{0.05}{\pu} for \SI{90}{\second} and
\SI{0.01}{\pu} for \SI{118}{\second}, after which the control stack has restored
the voltage and the \emph{settled} deviation is below \SI{0.01}{\pu}.

That timescale is exactly why the dead band matters. The TSO re-optimises every
\SI{180}{\second}; the whole excursion opens and closes \emph{inside one TSO
interval}. Whatever is done about it in that window is done by the DSO layer and
by the local droop --- which is the inter-OFO-step role the dead band governs.

\paragraph{Quasi-static screening cannot select the severe case.} A static
outage scan ranks the two machines in the opposite order to the dynamic plant:

\begin{center}
\begin{tabular}{lcc}
\toprule
worst-park $|\Delta V|$ [\pu]
  & static scan & RMS, droop disabled \\
  & (open loop, settled) & (closed loop, peak) \\
\midrule
trip gen 1 (\SI{650}{\mega\watt}) & 0.0854 & 0.1039 \\
trip gen 7 (\SI{830}{\mega\watt}) & \textbf{0.0122} & \textbf{0.2240} \\
\bottomrule
\end{tabular}
\end{center}

\noindent \textbf{The two columns are different quantities} and must not be read
as a ratio: the static scan is an unregulated power flow after the outage, while
the RMS column is the peak of a closed-loop transient. The comparison that
matters is the \emph{ordering}. Statically the larger machine looks like the
milder disturbance --- a distributed slack absorbs its \SI{830}{\mega\watt}
instantaneously and the post-outage power flow barely moves. Dynamically it
produces the deeper swing and is the only case reaching an undervoltage
violation. A severe case selected on the quasi-static plant would therefore have
been the wrong machine.

This is a concrete instance of the general point of \S\ref{sec:rev2}: the
controllers may reason on a reduced or quasi-static model, but the
\emph{evidence} must come from the plant that has the dynamics.

\subsection{The lower bound: what the droop must ignore}

Profile-driven drift is measured on the twins whose dead band is provably never
crossed ($\dz \geq 0.025$, four runs, $n=48$ TS and $3360$ DS inter-dispatch
windows). Admission is judged against the \emph{widest} twin's maximum, never
against a twin's own: a droop that successfully suppresses drift below its own
threshold would otherwise certify itself as inactive.

\begin{table}[htbp]
\centering
\begin{tabular}{lccc}
\toprule
 & p90 [\pu] & max [\pu] & windows exceeding $\dz$ \\
\midrule
TS parks & 0.00285 & 0.00304 & --- \\
DS parks & 0.00108 & 0.01101 & --- \\
\midrule
$\dz = 0.001$ & & & 28/48 TS,\; 499/3360 DS \\
$\dz = 0.005$ & & & \textbf{0/48} TS,\; 20/3360 DS (\SI{0.6}{\percent}) \\
$\dz = 0.01$  & & & 0/48 TS,\; 4/3360 DS (\SI{0.12}{\percent}) \\
$\dz = 0.025$ & & & 0/48 TS,\; 0/3360 DS \\
\bottomrule
\end{tabular}
\caption{Profile-driven drift and the resulting false-activation rate. The DS
distribution carries a $10\times$ tail (p90 0.0011, max 0.0110): the DS bound is
set by rare windows, not by typical operation.}
\label{tab:n1drift}
\end{table}

\subsection{The upper bound: what the droop must catch}

\begin{table}[htbp]
\centering
\begin{tabular}{lcccc}
\toprule
$\dz$ [\pu] & \multicolumn{2}{c}{compensation $1-\hat V/\hat V_{\text{no droop}}$}
            & \multicolumn{2}{c}{min.\ absolute $V$ [\pu]} \\
 & gen 1 & gen 7 & gen 1 & gen 7 \\
\midrule
0      & 0.78 & 0.55 & 1.006 & 0.935 \\
0.005  & 0.73 & 0.54 & 1.003 & 0.934 \\
0.01   & 0.67 & 0.54 & 0.996 & 0.933 \\
0.025  & 0.49 & 0.57 & 0.972 & 0.939 \\
0.05   & 0.28 & 0.52 & 0.966 & 0.916 \\
0.075  & 0.08 & 0.42 & 0.951 & \textbf{0.880} \\
0.15   & 0.00 & 0.00 & 0.937 & \textbf{0.816} \\
\bottomrule
\end{tabular}
\caption{Event rejection at the TS parks. With the droop disabled a gen-7 trip
drives a park terminal to \SI{0.816}{\pu}; the Q(V) layer prevents that
violation, but only while $\dz \lesssim 0.05$.}
\label{tab:n1reject}
\end{table}

\begin{figure}[htbp]
\centering
\includegraphics[width=0.78\textwidth]{fig_n1_rejection_gen7_20160105_0800.pdf}
\caption{Post-trip peak deviation and the residual still standing when the next
TSO dispatch acts, against $\dz$ (gen 7). The knee lies above
\SI{0.01}{\pu}.}
\label{fig:n1rejection}
\end{figure}

\begin{figure}[htbp]
\centering
\includegraphics[width=0.78\textwidth]{fig_n1_admissible_gen1_20160105_0800.pdf}
\caption{The two competing rates against $\dz$: false activation on profile
drift (left axis) and compensation effectiveness on the event (right axis). An
admissible $\dz$ lies where the first has fallen to zero and the second has not
yet decayed.}
\label{fig:n1admissible}
\end{figure}

\subsection{Choice of $\dz$}
\label{sec:n1choice}

Both bounds of \eqref{eq:detector} are measured at every window, and they do not
conflict: the admissible interval is non-empty everywhere. Within it, $\dz$ is a
quantile choice --- and replication is what fixes the quantile.

\paragraph{$\dz = \SI{0.005}{\pu}$ does not survive replication.} On the
reference window alone it looked ideal: TS false activation \emph{exactly} zero,
compensation within \SI{94}{\percent} of the ceiling. That zero is not a property
of the system but of that window. At the net-import window the TS parks drift
$2.4\times$ further (max \SI{0.0075}{\pu} against \SI{0.0031}{\pu}), and
$\dz=0.005$ then fires on \SI{13.3}{\percent} of TS inter-dispatch windows.

\begin{table}[htbp]
\centering
\begin{tabular}{lccc|ccc}
\toprule
& \multicolumn{3}{c|}{TS parks} & \multicolumn{3}{c}{DS parks} \\
$\dz$ [\pu] & $-117$ & $+409$ & $+1367$ & $-117$ & $+409$ & $+1367$ \\
\midrule
0.0025 & 0.450 & 0.167 & 0.083 & 0.060 & 0.030 & 0.046 \\
0.005  & \textbf{0.133} & 0.000 & 0.000 & 0.009 & 0.006 & 0.014 \\
0.01   & \textbf{0.000} & \textbf{0.000} & \textbf{0.000} & 0.000 & 0.001 & 0.005 \\
0.025  & 0.000 & 0.000 & 0.000 & 0.000 & 0.000 & 0.000 \\
\bottomrule
\end{tabular}
\caption{False-activation rate --- the fraction of inter-dispatch windows whose
open-loop drift exceeds $\dz$ --- at three operating windows labelled by net
infeed [\si{\mega\watt}]. Zero at the TS parks everywhere requires
$\dz \geq \SI{0.01}{\pu}$; zero at both levels requires $\dz \geq
\SI{0.0125}{\pu}$.}
\label{tab:n1fa}
\end{table}

\textbf{Recommendation: $\dz \approx \SI{0.01}{\pu}$}, on three grounds:

\begin{enumerate}
  \item \textbf{It is silent on profiles at every window tested} --- zero TS
    activation at all three, DS at most \SI{0.5}{\percent}
    (Table~\ref{tab:n1fa}). No smaller value has this property.
  \item \textbf{It lies inside the flat region of the rejection curve.}
    Compensation is flat up to $\dz \approx 0.025$ at every window and only then
    decays: at $\dz=0.01$ it is \numrange{0.54}{0.70} for the severe outage and
    \numrange{0.64}{0.67} for the milder one.
  \item \textbf{The cost of moving from \SI{0.005}{\pu} is small} --- about
    $0.05$ of compensation --- and it buys a property that holds at every window
    rather than at one.
\end{enumerate}

\noindent $\dz = \SI{0.0125}{\pu}$ additionally gives zero false activation at
the \emph{DS} parks, at slightly more compensation cost; the choice between them
is the risk preference over a \SI{0.5}{\percent} DS activation rate.

\paragraph{Why a quantile and not the maximum.} Anchoring to the drift
\emph{maximum} would be anchoring to the observation horizon: the same window
sampled over \SI{300}{\second} yields a DS maximum of \SI{0.0014}{\pu} against
\SI{0.0110}{\pu} over \SI{600}{\second}. A quantile is stable under that choice,
which is why Table~\ref{tab:n1fa} reports rates rather than extrema. Stated in
those terms:

\begin{quote}
At $\dz = \SI{0.01}{\pu}$ the local layer ignores \emph{all} profile-driven
inter-dispatch variation at the TS parks and \SI{99.5}{\percent} at the DS parks,
across net infeed from $-117$ to $+\SI{1367}{\mega\watt}$, while retaining
\numrange{0.54}{0.70} of the event compensation achievable at any dead band.
\end{quote}

\paragraph{What reproduces, and what does not.} The DS drift maximum is stable
(\num{0.0110}, \num{0.0089}, \num{0.0111}\,\pu, a \SI{24}{\percent} spread) and
so is the shape of the compensation curve. The TS drift maximum is \emph{not}
(\num{0.0075} against $\approx\!\num{0.0031}$, a \SI{147}{\percent} spread), and
that single fact is what moves the recommendation. Three candidate mechanisms
were tested and all refuted: DER reactive \emph{capability} is identical
(\SI{901.2}{\Mvar}) at all three windows with $P/S_n$ between \num{0.38} and
\num{0.60}, far above the \num{0.2} knee of the VDE diagram; $Q$ never
approaches its limit (peak utilisation \SI{12}{\percent}, no sample above
\SI{95}{\percent} of capability); and OLTC tap steps are both too rare
(\numrange{1}{3} per \SI{600}{\second} run) and associated with \emph{lower}
drift, not higher. The cause is unresolved and is stated as such.

\paragraph{One dead band, not two.} The TS admissible set
contains the DS set at this operating
point: the DS parks bind and the TS parks have margin. Separate $\dz_\mathrm{TS}$
and $\dz_\mathrm{DS}$ are therefore not required by this evidence. The per-level
machinery exists and a $\dz_\mathrm{TS}\!\times\!\dz_\mathrm{DS}$ sweep is the
natural follow-up, but the null result stands until it is run.

%======================================================================
\section{Threats to validity}

\begin{enumerate}
  \item \textbf{No measurement noise} (\S\ref{sec:excitation}). The mechanism a
    dead zone classically exists to reject is absent, so every $\dz^\star$ here
    is expected to be a lower bound for a noisy plant. This is the single change
    most likely to move the answer.
  \item \textbf{Equation~\eqref{eq:threshold} is verified at
    $|\Delta V| = 0.0026$ and \SI{0.0069}{\pu}.} Applying it at
    \SI{0.01}{\pu} extrapolates by about $1.5\times$. The extrapolation is
    modest and the relation is linear by construction, but it is unverified at
    that amplitude, and amplitudes large enough to test it directly destabilise
    the cascade.
  \item \textbf{One operating point for the threshold.} Whether $k\approx1$
    holds across loading is untested; the profiled windows validate the
    \emph{consequences} of a choice, not the threshold itself. Repeating
    \S\ref{sec:threshold} at a second window would settle it.
  \item \textbf{Purposive windows, $n=5$.} No frequency-weighted or
    annual-average statement is supportable. The set is also asymmetric about
    zero for the structural reason of \S\ref{sec:windows}.
  \item \textbf{Contingencies: addressed, at one window
    (\S\ref{sec:n1}).} A trip changes the topology, and therefore the
    controllers' cached sensitivities, whereas a load step changes only the
    operating point. Sections \ref{sec:sweep}--\ref{sec:rule} rest on load
    steps alone; \S\ref{sec:n1} adds N-1 evidence and reaches the same
    $\dz=\SI{0.005}{\pu}$ by an independent route. That evidence is
    \textbf{provisional}: one operating window, two machines.
  \item \textbf{The N-1 lower bound rests on a tail.} The DS false-activation
    rate at $\dz=\SI{0.005}{\pu}$ is \SI{0.6}{\percent} --- 20 windows out of
    3360. A bound set by twenty samples is sensitive to which windows are drawn,
    which is the strongest single argument for replicating \S\ref{sec:n1} across
    operating points before quoting the interval as general.
  \item \textbf{Two machines, not a contingency list.} \S\ref{sec:n1} trips two
    of six generators; \texttt{gen 9} is excluded because the static screening
    diverges on it. No N-1 \emph{completeness} claim is made, and line outages
    --- which the RMS adapter does not yet deliver --- are untested.
  \item \textbf{Residual TSO reduction error} of \SIrange{0.012}{0.018}{\pu}
    (\S\ref{sec:rev2}). The sensitivities are linearised once, and measurements
    come from the combined model, so an offset in a smooth regime is tolerable
    --- unlike the pre-correction DSO error, which placed the linearisation on
    the wrong branch of the $P$--$V$ curve.
  \item \textbf{The multi-equilibrium behaviour is characterised only
    qualitatively.} An ensemble over perturbed initial conditions per cell would
    map the basins directly and is the natural next experiment.
\end{enumerate}

\end{document}
